\documentclass{article}

\usepackage{course}
\usepackage{fullpage}
\usepackage{amsmath}
\usepackage{verbatim}
\usepackage{alltt}
\usepackage{url}

\newcommand{\code}[1]{{\tt #1}}
\newcommand{\class}[1]{{\tt #1}}

\def\handout{Lab 9}
\def\coursenum{142}
\def\coursetitle{Program design and methodology}
\def\courseterm{Spring}
\def\courseyear{2026}
\parskip 6pt
\parindent 0pt

\begin{document}

In this lab, we're returning to the fork-join solution technique, but
using it for operations other than sum.

\subsection*{Other Reductions using Fork-Join}

Begin by downloading the given code ({\tt lab9given.zip}).
Unzip this and open the {\tt lab9-reductions} folder it contains.
The only code this has initially is {\tt SumArray.java}, the same
program we worked with last week to sum the contents of an array.
(I did remove the restriction that all the array values be between 0
and 9.)

Today, we'll look at changing this code to perform different
operations.
All the operations will have a similar structure; they are all called
{\em reductions} because they reduce an array of input values to a
single result.

For the first modification, we will find the minimum.
You can either modify {\tt SumArray} or make a copy of it with a more
descriptive name.
Either way, change the compute function so that it returns the minimum
value rather than the sum.
This requires small changes to both the base case (find the minimum of
less than 100 values) and the recursive one (return the minimum of
{\tt left} and {\tt right} rather than their sum).
You will also want to change the serial computation in {\tt main} so
that you can check your answer.

Once you have that working, change the program again so that it
returns both the minimum and the index where the minimum value occurs.
To return multiple values, you'll need to create a class to store them
both.
Then use this class as the type of {\tt RecursiveTask} your code
extends (modify at the top of the {\tt SumArray} class) and as the
return value of {\tt compute}.
Then, you'll again modify the base case and recursion step of
{\tt compute} itself, as well as the serial implementation.

Next, we'll implement the reduction of finding the $k^\text{th}$
smallest value.
($k=1$ means finding the minimum, $k=2$ means finding the second
smallest, $k=3$ means finding the 3rd smallest, etc.)
Begin by implementing the serial version of this, possibly starting
specifically with $k=2$ (finding the second smallest value).
This is a generalization of finding the smallest value.
The algorithm for that was to traverse the array with a loop while
keeping track of the smallest value seen so far.
To find the $k^\text{th}$ smallest, use an array to keep track of the
$k$ smallest values seen so far (in order).
You can start by sorting the first $k$ values (the function
{\tt Arrays.sort} could be useful).
Then traverse the array, updating this array so that it always has the
$k$ smallest values so far.

Once you've got the serial code for $k^\text{th}$ smallest working,
implement it in the fork-join framework.
The object returned by {\tt compute} will now need to contain the
array with the $k$ smallest.
The base case will use your serial solution.
To combine {\tt left} and {\tt right} in the recursion step, you'll
want to do something like the merge operation in mergesort; the
smallest value is one of the smallest values in the left and right
solutions, the next smallest is either the other smallest value or the
second smallest in the array you took the overall smallest from, etc.

If you have more time, I have a challenge problem that goes beyond
reductions (but is related to them).
The goal is to implement {\em prefix sums}:
create a new array whose value at index $i$ is the sum of all the
values in the original array $\leq i$.
For example, if the input is
\begin{center}
3 1 4 2 5 3 1 3 4
\end{center}
then the array of prefix sums is
\begin{center}
3 4 8 10 15 18 19 22 26
\end{center}

To implement this in parallel, you'll want to go back to the original
{\tt SumArray} except that you should shrink the array and recursion
threshold so you can see what is happening.
I suggest you also use \% to restrict the array values to a small
range.
Then modify the code to take a second array; the original array is the
input and the new array will be the output.
You'll need 2 passes.
The first pass is essentially the existing {\tt SumArray} except that
it should write it's sum to the last cell of its part of the array in
addition to returning it.

For the second pass, you'll need a different extension of the
{\tt RecursiveTask} class with the same recursion structure as your
modified {\tt SumArray}. 
Tasks of this new class will take the sum of all the values before its
part of the array as an additional argument.
With this extra information, its base case will be able to fill in the second array.
For the recursion step, you'll need to read the array values written
in the first pass to calculate all the appropriate arguments.

\end{document}   
